\section{Discussion}
\label{sec:discussion}

\subsection{Synthesis of Findings}

The five experiments converge on a consistent picture: ink information in Scroll~1 CT data is confined to a narrow depth band around layers 30-34. The position sweep identifies this as the optimal window, with recall halving as the window shifts only two layers deeper. The leave-one-out analysis confirms that all five layers within this band contribute, with layers~32 and~31 carrying the strongest signal. Cross-segment evaluation shows the fixed window generalizes to 75\% of tested segments; the 25\% failure rate is attributable to depth offset or label misalignment data-level issues, not architectural limitations.

\subsection{Why Five Layers Suffice}

The GP Winner uses 26 input layers and solves depth selection implicitly through temporal attention, learning to suppress uninformative Z-layers. Our approach solves the same problem externally by hard-coding only the physically relevant layers. Both strategies converge on the same depth region; ours simply makes the selection explicit. The practical result is a $5.2\times$ reduction in input data with no loss of quality our 80-epoch UNet2D surpasses the GP Winner's $F_{0.5}$ by $25\%$. The ablation further shows that extending the window to $N = 9$ \emph{degrades} performance, consistent with the known dilution effect of uninformative input channels in multi-channel segmentation~\cite{roth2014new}.

\subsection{Architecture and Training}

Both architectures show identical qualitative behavior across all experiments: optimal at $N = 5$, $s = 30$, and peak test $F_{0.5}$ at 80 epochs followed by overfitting. This consistency supports the interpretation that findings reflect physical data properties rather than model-specific biases. The UNet2D consistently outperforms our TimeSformer at equal training budgets: with 5 input layers, the depth dimension is too shallow for divided space-time attention to add value over channel-wise convolution. Temporal attention assumes strong inter-frame dynamics, which adjacent CT depth layers lack. The validation-test divergence further demonstrates that a single validation segment is an unreliable checkpoint criterion on a small pseudo-label corpus.

\subsection{Practical Implications}

Using 5 layers instead of 26 reduces per-segment I/O by $5.2\times$, making systematic full-scroll scanning practical within academic compute budgets. The critical caveat is per-segment depth profiling: the 25\% failure rate in Experiment~D confirms that a static window fails when the true surface deviates from the assumed position. Computing the mean-intensity depth profile before inference and shifting the window to the intensity peak is a low-cost safeguard against this failure mode.

\subsection{Recommendations}

\begin{enumerate}
    \item \textbf{Use $N = 5$ layers, $s = 30$} (layers 30--34) as the default for Scroll~1 at 7.91~\textmu m resolution.
    \item \textbf{Clip intensities at $\tau = 200$} to suppress CT metal artifacts before training and inference.
    \item \textbf{Train for 80 epochs.} Both architectures overfit beyond this on a 4-segment pseudo-label corpus, with test $F_{0.5}$ declining $28\%$ from 80 to 100 epochs.
    \item \textbf{Profile each segment's depth curve} before inference. Shift the window if the intensity peak falls outside layers 30-34.
    \item \textbf{Prefer UNet2D over TimeSformer} at $N = 5$: higher $F_{0.5}$ (80ep: 0.474 vs.\ 0.403), simpler training, and no quadratic attention overhead.
\end{enumerate}

\section{Limitations}
\label{sec:limitations}

\paragraph{Pseudo-label quality.} Ground truth consists of GP model predictions rather than expert annotations. Training on pseudo-labels propagates and may amplify teacher errors through distillation~\cite{hinton2015distilling,lee2013pseudo}; the performance ceiling with genuine manual labels is unknown and likely higher.

\paragraph{Incomplete ablation.} The $N = 7$ and $N = 9$ runs were terminated after 7-9 epochs by SLURM time limits. While all configurations peaked at epochs 4-6, the absence of full training runs limits firm conclusions about the optimal $N$.

\paragraph{Small corpus, single scroll.} With 4 training segments from Scroll~1 only, findings may not transfer to scrolls with different scanning parameters or ink compositions~\cite{bukreeva2016virtual}. The validation-test divergence is a direct consequence of this small corpus.

\paragraph{Static Z-window.} A global fixed window cannot account for local surface curvature within a segment. The 25\% failure rate in Experiment~D reflects this limitation directly.

\section{Future Work}
\label{sec:future_work}

The most impactful extension is \emph{adaptive depth selection}: computing the mean-intensity depth profile per tile and centering the window on the detected surface peak, directly resolving the depth-offset failure mode. A learned variant could use channel attention to make Z-selection differentiable end-to-end.

Expanding the training corpus to 10-20 segments with iterative pseudo-label refinement following the GP Winner's 15-round approach~\cite{nader2024gp} would address the small-corpus limitations identified here. Cross-scroll validation on Scroll~5 would test whether the layer-30--34 window generalizes beyond Scroll~1.

\section{Conclusion}
\label{sec:conclusion}

This work presents the systematic study of Z-layer selection for ink detection in Herculaneum papyri CT data. We show that a 5-layer window centered at the papyrus surface (layers 30-34, ${\sim}40$~\textmu m) captures the essential ink signal. Leave-one-out sensitivity confirms that layers~32 and~31 carry the strongest signal, yet all five layers contribute and none is dispensable. A UNet2D trained for 80 epochs on pseudo-labels achieves $F_{0.5} = 0.474$ surpassing the competition-winning TimeSformer by $25\%$ using $5.2{\times}$ fewer input layers because focused depth selection raises signal-to-noise ratio where increased model complexity cannot. Training beyond 80 epochs on a 4-segment corpus induces overfitting that single-segment validation cannot detect, making epoch budget a critical and often overlooked design choice. Cross-segment evaluation demonstrates that generalization is achievable but brittle: six of eight tested segments succeed, while two fail due to depth offset or label misalignment data-level issues directly addressable through per-segment depth profiling. By translating competition code into an experimentally validated baseline, this work provides concrete, reproducible Z-layer configuration guidelines for future ink detection research.
